% Bagian: first-order-logic
% Bab: natural-deduction
% Bagian: proving-things

\documentclass[../../../include/open-logic-section]{subfiles}

\begin{document}

\iftag{FOL}
      {\olfileid[id]{fol}{ntd}{pro}}
      {\olfileid[id]{pl}{ntd}{pro}}

\olsection{Contoh \usetoken{P}{derivation}}

\begin{ex}
Mari kita berikan !!a{derivation} atas !!{sentence} $(!A \land !B) \lif !A$.

Kita mulai dengan menuliskan kesimpulan yang diinginkan di bagian
bawah !!{derivation}.
\begin{prooftree}
\AxiomC{}
\UnaryInfC{$(!A\land !B) \lif !A$}
\end{prooftree}

Selanjutnya, kita perlu menentukan jenis inferensi yang dapat
menghasilkan !!a{sentence} berbentuk demikian. !!^{main operator} dari
kesimpulan itu adalah $\lif$, sehingga kita akan mencoba mencapai
kesimpulan tersebut dengan aturan \Intro{\lif}. Sebaiknya kita
menuliskan asumsi-asumsi yang digunakan dan memberi label pada
aturan-aturan inferensi sambil melanjutkan, agar pada akhir pembuktian
mudah terlihat apakah semua asumsi telah !!{discharged}.
\begin{prooftree}
\AxiomC{$\Discharge{!A \land !B}{1}$}
\DeduceC{$!A$}
\DischargeRule{\Intro{\lif}}{1}
\UnaryInfC{$(!A\land !B) \lif !A$}
\end{prooftree}

Sekarang kita perlu melengkapi langkah-langkah dari asumsi $!A \land
!B$ menuju $!A$. Karena kita hanya menghadapi satu penghubung, yakni
$\land$, kita harus memakai aturan eliminasi $\land$. Hasilnya ialah
pembuktian berikut:
\begin{prooftree}
\AxiomC{$\Discharge{!A \land !B}{1}$}
\RightLabel{\Elim{\land}}
\UnaryInfC{$!A$}
\DischargeRule{\Intro{\lif}}{1}
\UnaryInfC{$(!A\land !B) \lif !A$}
\end{prooftree}
Sekarang kita memiliki !!a{derivation} yang benar atas $(!A \land
!B) \lif !A$.
\end{ex}

\begin{ex}
Sekarang mari kita berikan !!a{derivation} atas $(\lnot !A \lor !B)
\lif (!A \lif !B)$.

Kita mulai dengan menuliskan kesimpulan yang diinginkan di bagian
bawah !!{derivation}.
\begin{prooftree}
\AxiomC{}
\UnaryInfC{$(\lnot !A \lor !B) \lif (!A \lif !B)$}
\end{prooftree}
Untuk menemukan aturan logis yang dapat memberikan kesimpulan ini,
kita melihat penghubung-penghubung logis di dalamnya: $\lnot$, $\lor$,
dan $\lif$. Untuk sementara, kita hanya memedulikan kemunculan pertama
$\lif$ karena itulah !!{main operator} dari !!{sentence} yang menjadi
kesimpulan !!{derivation}; sedangkan $\lnot$, $\lor$, dan kemunculan
kedua $\lif$ berada dalam cakupan penghubung lain, sehingga akan kita
tangani kemudian. Oleh karena itu, kita mulai dengan aturan
\Intro{\lif}. Penerapan yang benar harus berbentuk seperti ini:
\begin{prooftree}
\AxiomC{$\Discharge{\lnot !A \lor !B}{1}$}
\DeduceC{$!A \lif !B$}
\DischargeRule{\Intro{\lif}}{1}
\UnaryInfC{$(\lnot !A \lor !B) \lif (!A \lif !B)$}
\end{prooftree}
Untuk melanjutkan, kini ada dua kemungkinan. Kita dapat terus bekerja
dari bawah ke atas dan mencari penerapan lain dari aturan
\Intro{\lif}, atau bekerja dari atas ke bawah dan menerapkan aturan
\Elim{\lor}. Mari kita pilih yang terakhir. Kita akan menggunakan
asumsi $\lnot !A \lor !B$ sebagai premis paling kiri dari \Elim{\lor}.
Agar penerapan \Elim{\lor} benar, kedua premis lainnya harus identik
dengan kesimpulan $!A \lif !B$, tetapi masing-masing dapat diturunkan
dari asumsi lain, yaitu salah satu dari dua disjung $\lnot !A \lor !B$.
Jadi, !!{derivation} kita akan tampak seperti ini:
\begin{prooftree}
\AxiomC{$\Discharge{\lnot !A \lor !B}{1}$}
\AxiomC{$\Discharge{\lnot !A}{2}$}
\DeduceC{$!A \lif !B$}
\AxiomC{$\Discharge{!B}{2}$}
\DeduceC{$!A \lif !B$}
\DischargeRule{\Elim{\lor}}{2}
\TrinaryInfC{$!A \lif !B$}
\DischargeRule{\Intro{\lif}}{1}
\UnaryInfC{$(\lnot !A \lor !B) \lif (!A \lif !B)$}
\end{prooftree}

Pada masing-masing dari dua cabang di sebelah kanan, kita hendak
!!{derive} $!A \lif !B$, yang paling baik dilakukan dengan
\Intro{\lif}.
\begin{prooftree}
\AxiomC{$\Discharge{\lnot !A \lor !B}{1}$}
\AxiomC{$\Discharge{\lnot !A}{2}, \Discharge{!A}{3}$}
\DeduceC{$!B$}
\DischargeRule{\Intro{\lif}}{3}
\UnaryInfC{$!A \lif !B$}
\AxiomC{$\Discharge{!B}{2}, \Discharge{!A}{4}$}
\DeduceC{$!B$}
\DischargeRule{\Intro{\lif}}{4}
\UnaryInfC{$!A \lif !B$}
\DischargeRule{\Elim{\lor}}{2}
\TrinaryInfC{$!A \lif !B$}
\DischargeRule{\Intro{\lif}}{1}
\UnaryInfC{$(\lnot !A \lor !B) \lif (!A \lif !B)$}
\end{prooftree}

Untuk dua bagian !!{derivation} yang masih hilang, kita memerlukan
!!{derivation}s atas $!B$ dari $\lnot !A$ dan $!A$ pada cabang tengah,
serta dari $!A$ dan $!B$ pada cabang kanan. Mari kita tangani yang
pertama lebih dahulu. $\lnot !A$ dan $!A$ merupakan dua premis
\Elim{\lnot}:
\begin{prooftree}
\AxiomC{$\Discharge{\lnot !A}{2}$}
\AxiomC{$\Discharge{!A}{3}$}
\RightLabel{\Elim{\lnot}}
\BinaryInfC{$\lfalse$}
\DeduceC{$!B$}
\end{prooftree}
Dengan menggunakan \FalseInt, kita dapat memperoleh $!B$ sebagai
kesimpulan dan menyelesaikan cabang tersebut.
\begin{prooftree}
\AxiomC{$\Discharge{\lnot !A \lor !B}{1}$}
\AxiomC{$\Discharge{\lnot !A}{2}$}
\AxiomC{$\Discharge{!A}{3}$}
\RightLabel{\Elim{\lnot}}
\BinaryInfC{$\lfalse$}
\RightLabel{\FalseInt}
\UnaryInfC{$!B$}
\DischargeRule{\Intro{\lif}}{3}
\UnaryInfC{$!A \lif !B$}
\AxiomC{$\Discharge{!B}{2}, \Discharge{!A}{4}$}
\DeduceC{$!B$}
\DischargeRule{\Intro{\lif}}{4}
\UnaryInfC{$!A \lif !B$}
\DischargeRule{\Elim{\lor}}{2}
\TrinaryInfC{$!A \lif !B$}
\DischargeRule{\Intro{\lif}}{1}
\UnaryInfC{$(\lnot !A \lor !B) \lif (!A \lif !B)$}
\end{prooftree}

Sekarang mari kita periksa cabang paling kanan. Di sini penting untuk
menyadari bahwa definisi !!{derivation} \emph{memungkinkan asumsi untuk
  dilepaskan}, tetapi \emph{tidak mewajibkan} pelepasan itu. Dengan kata
lain, jika kita dapat menurunkan $!B$ dari salah satu asumsi $!A$ dan
$!B$ tanpa menggunakan yang lain, hal itu diperbolehkan. !!^{derive}
$!B$ dari~$!B$ pun sepele: $!B$ sendiri merupakan !!a{derivation}
semacam itu dan tidak diperlukan inferensi apa pun. Jadi, kita cukup
menghapus asumsi~$!A$.
\begin{prooftree}
\AxiomC{$\Discharge{\lnot !A \lor !B}{1}$}
\AxiomC{$\Discharge{\lnot !A}{2}$}
\AxiomC{$\Discharge{!A}{3}$}
\RightLabel{\Elim{\lnot}}
\BinaryInfC{$\lfalse$}
\RightLabel{\FalseInt}
\UnaryInfC{$!B$}
\DischargeRule{\Intro{\lif}}{3}
\UnaryInfC{$!A \lif !B$}
\AxiomC{$\Discharge{!B}{2}$}
\RightLabel{\Intro{\lif}}
\UnaryInfC{$!A \lif !B$}
\DischargeRule{\Elim{\lor}}{2}
\TrinaryInfC{$!A \lif !B$}
\DischargeRule{\Intro{\lif}}{1}
\UnaryInfC{$(\lnot !A \lor !B) \lif (!A \lif !B)$}
\end{prooftree}
Perhatikan bahwa dalam !!{derivation} yang sudah selesai, inferensi
\Intro{\lif} paling kanan sebenarnya tidak melepaskan asumsi apa pun.
\end{ex}

\begin{ex}
Sejauh ini kita belum memerlukan aturan \FalseCl{}. Aturan ini istimewa
karena memungkinkan kita melepaskan asumsi yang bukan sub!!{formula} dari
kesimpulan aturan. Aturan ini berkaitan erat dengan aturan \FalseInt{}.
Faktanya, aturan \FalseInt{} merupakan kasus khusus dari aturan
\FalseCl{}---terdapat logika yang disebut ``logika intuisionistik'' yang
hanya mengizinkan \FalseInt{}. Aturan \FalseCl{} merupakan jalan
terakhir ketika cara lain tidak berhasil. Sebagai contoh, andaikan kita
ingin !!{derive} $!A \lor \lnot !A$. Strategi kita biasanya ialah
mencoba !!{derive} $!A \lor \lnot !A$ dengan menggunakan
$\Intro{\lor}$. Namun, ini mengharuskan kita !!{derive} $!A$ atau
$\lnot !A$ tanpa asumsi, dan hal itu tidak dapat dilakukan. Di sinilah
\FalseCl{} menolong!
\begin{prooftree}
  \AxiomC{$\Discharge{\lnot(!A \lor \lnot !A)}{1}$}
  \DeduceC{$\lfalse$}
  \DischargeRule{\FalseCl}{1}
  \UnaryInfC{$!A \lor \lnot !A$}
\end{prooftree}
Sekarang kita mencari !!a{derivation} atas $\lfalse$ dari $\lnot(!A
\lor \lnot !A)$. Karena $\lfalse$ merupakan kesimpulan
$\Elim{\lnot}$, kita dapat mencoba aturan tersebut:
\begin{prooftree}
  \AxiomC{$\Discharge{\lnot(!A \lor \lnot !A)}{1}$}
  \DeduceC{$\lnot !A$}
  \AxiomC{$\Discharge{\lnot(!A \lor \lnot !A)}{1}$}
  \DeduceC{$!A$}
  \RightLabel{\Elim{\lnot}}
  \BinaryInfC{$\lfalse$}
  \DischargeRule{\FalseCl}{1}
  \UnaryInfC{$!A \lor \lnot !A$}
\end{prooftree}
Strategi kita untuk menemukan !!a{derivation} atas~$\lnot !A$
memerlukan penerapan~$\Intro{\lnot}$:
\begin{prooftree}
  \AxiomC{$\Discharge{\lnot(!A \lor \lnot !A)}{1}, \Discharge{!A}{2}$}
  \DeduceC{$\lfalse$}
  \DischargeRule{\Intro{\lnot}}{2}
  \UnaryInfC{$\lnot !A$}
  \AxiomC{$\Discharge{\lnot(!A \lor \lnot !A)}{1}$}
  \DeduceC{$!A$}
  \RightLabel{\Elim{\lnot}}
  \BinaryInfC{$\lfalse$}
  \DischargeRule{\FalseCl}{1}
  \UnaryInfC{$!A \lor \lnot !A$}
\end{prooftree}
Di sini, kita dapat memperoleh $\lfalse$ dengan mudah melalui penerapan
$\Elim{\lnot}$ pada asumsi $\lnot(!A \lor \lnot !A)$ dan $!A \lor
\lnot !A$, yang mengikuti dari asumsi baru kita $!A$ berdasarkan
~$\Intro{\lor}$:
\begin{prooftree}
  \AxiomC{$\Discharge{\lnot(!A \lor \lnot !A)}{1}$}
  \AxiomC{$\Discharge{!A}{2}$}
  \RightLabel{\Intro{\lor}}
  \UnaryInfC{$!A \lor \lnot !A$}
  \RightLabel{\Elim{\lnot}}
  \BinaryInfC{$\lfalse$}
  \DischargeRule{\Intro{\lnot}}{2}
  \UnaryInfC{$\lnot !A$}
  \AxiomC{$\Discharge{\lnot(!A \lor \lnot !A)}{1}$}
  \DeduceC{$!A$}
  \RightLabel{\Elim{\lnot}}
  \BinaryInfC{$\lfalse$}
  \DischargeRule{\FalseCl}{1}
  \UnaryInfC{$!A \lor \lnot !A$}
\end{prooftree}
Di sisi kanan kita menggunakan strategi yang sama, hanya saja kita
memperoleh $!A$ dengan~\FalseCl:
\begin{prooftree}
  \AxiomC{$\Discharge{\lnot(!A \lor \lnot !A)}{1}$}
  \AxiomC{$\Discharge{!A}{2}$}
  \RightLabel{\Intro{\lor}}
  \UnaryInfC{$!A \lor \lnot !A$}
  \RightLabel{\Elim{\lnot}}
  \BinaryInfC{$\lfalse$}
  \DischargeRule{\Intro{\lnot}}{2}
  \UnaryInfC{$\lnot !A$}
  \AxiomC{$\Discharge{\lnot(!A \lor \lnot !A)}{1}$}
  \AxiomC{$\Discharge{\lnot !A}{3}$}
  \RightLabel{\Intro{\lor}}
  \UnaryInfC{$!A \lor \lnot !A$}
  \RightLabel{\Elim{\lnot}}
  \BinaryInfC{$\lfalse$}
  \DischargeRule{\FalseCl}{3}
  \UnaryInfC{$!A$}
  \RightLabel{\Elim{\lnot}}
  \BinaryInfC{$\lfalse$}
  \DischargeRule{\FalseCl}{1}
  \UnaryInfC{$!A \lor \lnot !A$}
\end{prooftree}
\end{ex}

\begin{prob}
Berikan !!{derivation}s yang menunjukkan hal-hal berikut:
\begin{enumerate}
\item $!A \land (!B \land !C) \Proves (!A \land !B) \land !C$.
\item $!A \lor (!B \lor !C) \Proves (!A \lor !B) \lor !C$.
\item $!A \lif (!B \lif !C) \Proves !B \lif (!A \lif !C)$.
\item $!A \Proves \lnot\lnot !A$.
\end{enumerate}
\end{prob}

\begin{prob}
Berikan !!{derivation}s yang menunjukkan hal-hal berikut:
\begin{enumerate}
\item $(!A \lor !B) \lif !C \Proves !A \lif !C$.
\item $(!A \lif !C) \land (!B \lif !C) \Proves (!A \lor !B) \lif !C$.
\item $\Proves \lnot(!A \land \lnot !A)$.
\item $!B \lif !A \Proves \lnot !A \lif \lnot !B$.
\item $\Proves (!A \lif \lnot !A) \lif \lnot !A$.
\item $\Proves \lnot(!A \lif !B) \lif \lnot !B$.
\item $!A \lif !C \Proves \lnot (!A \land \lnot !C)$.
\item $!A \land \lnot !C \Proves \lnot (!A \lif !C)$.
\item $!A \lor !B, \lnot !B \Proves !A$.
\item $\lnot !A \lor \lnot !B \Proves \lnot(!A \land !B)$.
\item $\Proves (\lnot !A \land \lnot !B) \lif\lnot(!A \lor !B)$.
\item $\Proves \lnot(!A \lor !B) \lif (\lnot !A \land \lnot !B)$.
\end{enumerate}
\end{prob}

\begin{prob}
Berikan !!{derivation}s yang menunjukkan hal-hal berikut:
\begin{enumerate}
\item $\lnot(!A \lif !B) \Proves !A$.
\item $\lnot(!A \land !B) \Proves \lnot !A \lor \lnot !B$.
\item $!A \lif !B \Proves \lnot !A \lor !B$.
\item $\Proves \lnot \lnot !A \lif !A$.
\item $!A \lif !B, \lnot !A \lif !B \Proves !B$.
\item $(!A \land !B) \lif !C \Proves (!A \lif !C) \lor (!B \lif !C)$.
\item $(!A \lif !B) \lif !A \Proves !A$.
\item $\Proves (!A \lif !B) \lor (!B \lif !C)$.
\end{enumerate}
(Semua ini memerlukan aturan~$\FalseCl$.)
\end{prob}

\end{document}
